\section{Метод}\label{sec:method}

\begin{definition}[расстояние промаха]
След $\hat F_k$ представлен центрами клеток, в которые попадают точки истинных объектов класса $k$, перенесённые в систему планировщика; клетки за краем карты сохраняются. Для сцены и класса $k$
\begin{equation}
  s_k = \min\bigl\{ R_{\max},\ h\bigl(\hat F_k, A_\lambda(k)\bigr) \bigr\},
  \qquad h(X, Y) = \max_{x \in X}\, \min_{y \in Y} \|x - y\|,
  \label{eq:miss}
\end{equation}
где $h$ --- направленное расстояние Хаусдорфа. Если объектов класса в сцене нет, $s_k = 0$; если область пуста, $s_k = R_{\max}$.
\end{definition}

Промах --- расстояние от самой дальней точки истинного объекта до того, что нарисовано на карте (рис.~\ref{fig:example}). Пропущенный или перепутанный объект даёт большой промах, недорисованный --- промах на величину недорисованной части, сдвиг карты --- промах порядка сдвига, а поворот ещё и с плечом: дальние от робота объекты смещаются сильнее. Промах считается один на сцену: клетки, кадры и объекты одной сцены сильно зависимы, и считать их отдельными примерами нельзя (допущение~\ref{as:exch}).

\textbf{Калибровка.} Пусть $s^{(1)} \le \dots \le s^{(n)}$ --- упорядоченные промахи калибровочных сцен для класса $k$ и $m = \lceil (n + 1)(1 - \alpha) \rceil$. Запас
\begin{equation}
  \hat r_k = s^{(m)}, \quad \text{если } m \le n \text{ и } s^{(m)} < R_{\max}; \qquad \hat r_k = \infty \text{ иначе (отказ)}.
\end{equation}
При $n = 13$ и $\alpha = 0{,}1$ ранг $m = 13$, то есть запас равен наибольшему из калибровочных промахов. Калибровка ведётся отдельно по классам, чтобы ненадёжный класс не раздувал запас остальных.

\begin{proposition}\label{prop:valid}
При допущениях~\ref{as:exec} и~\ref{as:exch}
\[
  \mathbb P\bigl(\hat r_k = \infty \ \text{или}\ s_k(\Omega_{n+1}) \le \hat r_k\bigr) \ge \frac{m}{n + 1} \ge 1 - \alpha,
\]
и запас $\hat r_k$ удовлетворяет~\eqref{eq:goal} для путей по центрам клеток.
\end{proposition}
\begin{proof}
По обменяемости ранг $s_k(\Omega_{n+1})$ среди $n + 1$ промахов равномерно распределён; совпадения значений только помогают. Поэтому $\mathbb P(s_k(\Omega_{n+1}) \le s^{(m)}) \ge m / (n + 1)$~\cite{gentle}, а при отказе событие выполнено. Пусть $\hat r_k < \infty$ и $s_k \le \hat r_k$. Тогда для любой точки $y \in \hat F_k$ найдётся $a \in A_\lambda(k)$ с $\|y - a\| \le \hat r_k$, и для точки пути $x$ с $\operatorname{dist}(x, A_\lambda(k)) > d + \hat r_k$ по неравенству треугольника $\|x - y\| \ge \|x - a\| - \|a - y\| > d$.
\end{proof}

Гарантия не зависит от планировщика: от него зависит только, найдёт ли он путь в суженном свободном пространстве. Она относится к центрам клеток. Для непрерывного следа и непрерывного движения между центрами расстояние может оказаться меньше $d$ не больше чем на $\sqrt 2\,\Delta \approx 7$~см (по $\Delta/\sqrt 2$ на представление следа и на отрезок пути). Гарантия поклассовая: для миссии с несколькими запретными классами $\mathcal K_m$ по неравенству Буля все классы покрыты с вероятностью не ниже $1 - |\mathcal K_m|\,\alpha$. Она маргинальна: усреднена по калибровочным выборкам, и при $n = 13$ покрытие отдельной выборки распределено как $\mathrm{Beta}(13, 1)$ со средним 0,93 и 10-м перцентилем 0,84~\cite{gentle}.

\textbf{Сравниваемые калибровки.} Варианты отличаются областью и тем, по каким промахам считается запас:
\begin{itemize}[nosep]
  \item \emph{без калибровки}: область $A_{\lambda_0}(k)$, запас 0;
  \item \emph{только метки}: промах~\eqref{eq:miss} карты детектора при точных позах ($E_t \equiv I$), запас $\hat r^{\mathrm{lab}}_k$;
  \item \emph{только поза}: промах карты с эталонными метками, построенной по позам с дрейфом, запас $\hat r^{\mathrm{pose}}_k$;
  \item \emph{сумма}: $\hat r^{\mathrm{lab}}_k + \hat r^{\mathrm{pose}}_k$, оба на уровне $\alpha$;
  \item \emph{совместный} (наш метод): промах карты детектора с дрейфом, запас $\hat r^{\mathrm{joint}}_k$;
  \item \emph{только геометрия}: областью класса считаются все занятые клетки ($\lambda = 0$), запас --- квантиль промаха до них на карте с дрейфом. Это аналог раздувания карты занятости~\cite{pwc}: гарантия без семантики;
  \item \emph{множества меток} по~\cite{sundarsingh}: мера несоответствия сцены --- худшее $1 - p_j(k)$ по клеткам истинного следа ($\infty$, если клетка не занята), $\hat\eta$ --- её конформный квантиль; запретны клетки с $p_j(k) \ge 1 - \hat\eta$, запаса в метрах нет. У авторов максимум берётся по всем наблюдённым клеткам и всем классам, включая «свободно»; наше ограничение следом может только уменьшить меру.
\end{itemize}
Первые пять вариантов используют одну область $A_{\lambda_0}(k)$ с порогом $\lambda_0 = 0{,}05$, выбранным на отдельной сцене. Варианты без запаса считаются покрывшими сцену, если промах не больше одной клетки, --- это в их пользу.

\begin{proposition}\label{prop:sep}
Пусть для новой сцены
\begin{equation}
  s_k \le s^{\mathrm{lab}}_k + s^{\mathrm{pose}}_k,
  \label{eq:additive}
\end{equation}
где слагаемые --- промахи вариантов «только метки» и «только поза» той же сцены. Тогда $\mathbb P\bigl(s_k \le \hat r^{\mathrm{lab}}_k + \hat r^{\mathrm{pose}}_k\bigr) \ge 1 - 2\alpha$, а не $1 - \alpha$.
\end{proposition}
\begin{proof}
Каждый запас не покрывает свой промах с вероятностью не больше $\alpha$; по неравенству Буля оба покрывают с вероятностью не ниже $1 - 2\alpha$, и тогда по~\eqref{eq:additive} $s_k \le \hat r^{\mathrm{lab}}_k + \hat r^{\mathrm{pose}}_k$.
\end{proof}

Условие~\eqref{eq:additive} из геометрии не следует. По неравенству треугольника $s_k \le s^{\mathrm{lab}}_k + h(E_Q A^0, A)$, где $A^0$ --- область карты детектора при точных позах: второе слагаемое --- сдвиг \emph{предсказанной} области. Вариант «только поза» измеряет другое --- промах эталонной карты, как сделал бы инженер, отлаживающий локализацию отдельно от восприятия. На наших данных~\eqref{eq:additive} нарушается в части реализаций (разд.~\ref{sec:composition}). Распределение риска $\alpha/2 + \alpha/2$ дало бы $1 - \alpha$~\cite{calafiore2026}, но требует $n \ge 2/\alpha - 1$, то есть 19 калибровочных сцен при $\alpha = 0{,}1$. Где окажется покрытие суммы, решает неизвестная зависимость двух ошибок; совместный запас от неё не зависит.

\textbf{Миссии.} Робот радиуса 0,2~м едет от старта до цели. Запретная зона класса --- центры клеток на расстоянии не больше $d + \hat r_k$ от области класса; при отказе класса используется запасной сертификат «только геометрия». Путь ищется алгоритмом Дейкстры на 8-связной сетке. Он полон на графе сетки, поэтому «путь не найден» означает, что пути по свободным клеткам нет. Путь \emph{опасен}, если его центр клетки ближе $d$ к центру клетки истинного запретного объекта в системе планировщика. Путь \emph{сталкивается}, если проходит ближе $0{,}2 - \Delta$~м к истинному препятствию: карта препятствий не калибруется, и сдвиг на одну клетку столкновением не считается. Миссия \emph{успешна}, если путь найден, не опасен и не сталкивается. Кратчайший путь проверяет гарантию лишь вдоль одной линии, поэтому мы считаем и худший случай по всем планировщикам: есть ли клетка ближе $d$ к истинному объекту, достижимая от старта без входа в запретную зону. Тот же планировщик на истинной карте даёт оракул.

\textbf{Гипотезы.} Обозначим $\mathrm{Cov}(r)$ левую часть~\eqref{eq:goal}.
\begin{enumerate}[label=H\arabic*., nosep]
  \item Запас «только метки» даёт $\mathrm{Cov} < 1 - \alpha$, когда дрейф сравним с ошибкой меток.
  \item Запас «только поза» даёт $\mathrm{Cov} < 1 - \alpha$, пока дрейф меньше ошибки меток.
  \item Совместный запас не больше суммы: $\hat r^{\mathrm{joint}}_k \le \hat r^{\mathrm{lab}}_k + \hat r^{\mathrm{pose}}_k$.
  \item Метрика качества карты mIoU слабо связана с безопасностью миссий.
\end{enumerate}
Покрытие совместного варианта не гипотеза: оно следует из утверждения~\ref{prop:valid}, и его измерение проверяет реализацию.
